\documentclass[14pt,a4paper,oneside]{extbook}
\usepackage[utf8]{vietnam}
\usepackage{amsmath,amssymb,mathrsfs,fancyhdr,enumerate,multirow,makecell}
\usepackage{tkz-euclide,tikz-3dplot,tikz,tkz-tab}
\usetikzlibrary{arrows.meta,calc,intersections,angles,quotes,backgrounds,shapes.geometric,patterns,through,shadows,positioning}
\usepackage{yhmath}
\usepackage[top=0.7cm, bottom=0.7cm, left=1.0cm, right=1.0cm, headheight=16pt, headsep=3mm, footskip=5mm]{geometry}
\usepackage[hidelinks,unicode]{hyperref}
\usepackage{qrcode}
\usepackage[many]{tcolorbox}
\usepackage{fontawesome}

\pagestyle{fancy}
\fancyhf{}
\fancyhead[L]{\footnotesize\bfseries\color{teal!85!black}\faBook~SỔ TAY KIẾN THỨC TOÁN 4}
\fancyhead[R]{\footnotesize\bfseries\color{blue!85!black}Lớp Toán Thầy Lĩnh -- loptoanthaylinh.vn}
\fancyfoot[C]{\footnotesize\bfseries\color{gray!80!black}-- Sổ tay kiến thức: Bài 69. Ôn tập phân số --}
\renewcommand{\headrulewidth}{0.4pt}

\begin{document}

% --- BANNER TIÊU ĐỀ BÀI HỌC ---
\noindent
\begin{tcolorbox}[
  enhanced,
  colback=teal!8!white,
  colframe=teal!80!black,
  arc=3mm, boxrule=1pt,
  left=3mm, right=3mm, top=2mm, bottom=2mm,
  halign=center
]
  {\Large\bfseries\color{teal!90!black}\faBookmark\ BÀI 69: ÔN TẬP PHÂN SỐ}\\[1.5pt]
  {\footnotesize\itshape\color{black!75} Bản chất phân số, rút gọn tối giản, quy đồng mẫu số và so sánh phân số trên 1 trang A4}
\end{tcolorbox}

\vspace{-2.5mm}

\noindent
\begin{minipage}[t]{0.49\textwidth}\vspace{0pt}
\begin{tcolorbox}[
  enhanced,
  colback=blue!4!white,
  colframe=blue!75!black,
  arc=2mm, boxrule=0.8pt,
  left=2mm, right=2mm, top=1.8mm, bottom=1.8mm,
  title={\footnotesize\bfseries\color{white}\faCheckSquare\ 1. Tính chất \& Rút gọn phân số}
]
{\footnotesize
\begin{itemize}\setlength{\itemsep}{1pt}\setlength{\leftskip}{-2.5mm}
  \item \textbf{Tính chất cơ bản:}
  $$\dfrac{a}{b} = \dfrac{a \times m}{b \times m} = \dfrac{a : n}{b : n} \quad (m, n \ne 0)$$
  \item \textbf{Rút gọn về tối giản:} Chia cả tử và mẫu cho ước chung lớn hơn 1 cho đến khi tử và mẫu không cùng chia hết cho số nào khác 1.
  \item \textbf{Phân số với 1:} Tử $<$ Mẫu $\implies <1$; Tử $=$ Mẫu $\implies =1$; Tử $>$ Mẫu $\implies >1$.
\end{itemize}
}
\end{tcolorbox}
\end{minipage}
\hfill
\begin{minipage}[t]{0.49\textwidth}\vspace{0pt}
\begin{tcolorbox}[
  enhanced,
  colback=teal!4!white,
  colframe=teal!75!black,
  arc=2mm, boxrule=0.8pt,
  left=2mm, right=2mm, top=1.8mm, bottom=1.8mm,
  title={\footnotesize\bfseries\color{white}\faArrowsH\ 2. Quy đồng \& So sánh phân số}
]
{\footnotesize
\begin{itemize}\setlength{\itemsep}{1.5pt}\setlength{\leftskip}{-2.5mm}
  \item \textbf{Quy đồng mẫu số:} Chọn MSC nhỏ nhất (kiểm tra mẫu lớn chia hết mẫu bé trước khi nhân chéo).
  \item \textbf{So sánh phân số:}
  \begin{itemize}\setlength{\itemsep}{0.5pt}\setlength{\leftskip}{-2.5mm}
    \item Cùng mẫu: Tử lớn hơn $\to$ Phân số lớn hơn.
    \item Cùng tử: Mẫu BÉ hơn $\to$ Phân số LỚN hơn.
    \item Khác mẫu: Quy đồng rồi so sánh tử.
  \end{itemize}
\end{itemize}
}
\end{tcolorbox}
\end{minipage}

\vspace{-2.5mm}
% --- KHỐI 3: SƠ ĐỒ TƯ DUY TIKZ ---
\noindent
\begin{tcolorbox}[
  enhanced,
  colback=orange!4!white,
  colframe=orange!80!black,
  arc=2mm, boxrule=0.8pt,
  left=2mm, right=2mm, top=1.5mm, bottom=1.5mm,
  title={\footnotesize\bfseries\color{white}\faSitemap\ SƠ ĐỒ TƯ DUY: BẢN ĐỒ 4 TRỤ CỘT MẠCH KIẾN THỨC PHÂN SỐ}
]
\centering
\begin{tikzpicture}[
  box/.style={draw, rounded corners=2pt, align=center, font=\scriptsize\bfseries, inner sep=3pt, minimum width=3.4cm, minimum height=0.9cm},
  arr/.style={-{Stealth[length=2mm]}, thick}
]
  \node[draw=purple!80!black, fill=purple!15, rounded corners=3pt, font=\small\bfseries\color{purple!90!black}, inner sep=4pt, align=center] (core) at (5, 0.7) {TRỌNG TÂM\\PHÂN SỐ};

  \node[box, fill=blue!10, draw=blue!80!black] (k1) at (0.8, 1.4) {1. BẢN CHẤT\\Tử / Mẫu ($b > 0$)};
  \node[box, fill=teal!10, draw=teal!80!black] (k2) at (9.2, 1.4) {2. TỐI GIẢN\\Rút gọn hết cỡ};
  \node[box, fill=orange!12, draw=orange!80!black] (k3) at (0.8, 0.0) {3. QUY ĐỒNG\\Đưa về cùng mẫu};
  \node[box, fill=green!15, draw=green!70!black] (k4) at (9.2, 0.0) {4. SO SÁNH\\3 Ngã rẽ so sánh};

  \draw[arr, blue!80!black] (core) -- (k1);
  \draw[arr, teal!80!black] (core) -- (k2);
  \draw[arr, orange!80!black] (core) -- (k3);
  \draw[arr, green!70!black] (core) -- (k4);
\end{tikzpicture}
\end{tcolorbox}

\vspace{-2.5mm}
% --- KHỐI 4: VÍ DỤ & MẸO TRÁNH BẪY ---
\noindent
\begin{minipage}[t]{0.49\textwidth}\vspace{0pt}
\begin{tcolorbox}[
  enhanced,
  colback=green!4!white,
  colframe=green!60!black,
  arc=2mm, boxrule=0.8pt,
  left=2mm, right=2mm, top=1.8mm, bottom=1.8mm,
  title={\footnotesize\bfseries\color{white}\faCheckCircle\ Ví dụ thực hành chuẩn mực}
]
{\footnotesize
\textbf{1. Rút gọn:} $\dfrac{36}{48} = \dfrac{36 : 12}{48 : 12} = \mathbf{\dfrac{3}{4}}$ (tối giản).\\
\textbf{2. Quy đồng và so sánh:} So sánh $\dfrac{5}{6}$ và $\dfrac{7}{9}$:\\
$\text{MSC} = 18 \implies \dfrac{5}{6} = \dfrac{15}{18}; \quad \dfrac{7}{9} = \dfrac{14}{18}$.\\
Vì $\dfrac{15}{18} > \dfrac{14}{18} \implies \mathbf{\dfrac{5}{6} > \dfrac{7}{9}}$.
}
\end{tcolorbox}
\end{minipage}
\hfill
\begin{minipage}[t]{0.49\textwidth}\vspace{0pt}
\begin{tcolorbox}[
  enhanced,
  colback=red!4!white,
  colframe=red!70!black,
  arc=2mm, boxrule=0.8pt,
  left=2mm, right=2mm, top=1.8mm, bottom=1.8mm,
  title={\footnotesize\bfseries\color{white}\faExclamationTriangle\ Mẹo tránh bẫy trắc nghiệm}
]
{\footnotesize
\begin{itemize}\setlength{\itemsep}{1.5pt}\setlength{\leftskip}{-2.5mm}
  \item \textbf{Bẫy so sánh cùng tử số:}\\
  $\frac{5}{8}$ và $\frac{5}{11} \implies \frac{5}{8} > \frac{5}{11}$ (vì $8 < 11$). Cùng tử thì mẫu bé hơn phân số sẽ lớn hơn!
  \item \textbf{Bẫy phân số đảo ngược:}\\
  Số tự nhiên 1 có phân số đảo ngược là chính nó ($1 = \frac{1}{1}$). Số 0 \textbf{không có} phân số đảo ngược!
\end{itemize}
}
\end{tcolorbox}
\end{minipage}

\end{document}
